\documentclass[12pt, twoside]{article}
\input{../preamble}

\fancyhead[LO]{La Scuola d'Italia / Dr.\ Huson / IB Math AA SL \\* 5 May 2026}
\fancyhead[RO]{\fbox{\begin{minipage}{4cm}\footnotesize First \& Last name:\\[0.5cm]\end{minipage}}}
\fancyfoot[C]{\thepage}

\begin{document}
\subsubsection*{11.16 Classwork: Distributions}

Problems 1 and 2 are no calculator. Problems 3, 4, 5 calculator allowed.

\begin{enumerate}[itemsep=0.3cm]

%--- Problem 1: 23N.1.SL.TZ1.6 (P1 SL, no calc) ---
\item The binomial expansion of $(1 + kx)^{n}$ is given by
      \[ 1 + \tfrac{9x}{2} + 15k^{2}x^{2} + \cdots + k^{n}x^{n}, \]
      where $n \in \mathbb{Z}^{+}$ and $k \in \mathbb{Q}$.

      Find the value of $n$ and the value of $k$. \hfill [6]

\begin{tikzpicture}
  \draw (-0.4,0) rectangle (15,7);
  \foreach \y in {1,...,6} \draw[dotted] (0,\y)--(14.5,\y);
\end{tikzpicture}

%--- Problem 2: 24N.2.SL.TZ2.2 (P2 SL, calc allowed; doable by hand) ---
\item Find the coefficient of $x^{6}$ in the expansion of $(2x - 5)^{9}$. \hfill [4]

\begin{tikzpicture}
  \draw (-0.4,0) rectangle (15,7);
  \foreach \y in {1,...,6} \draw[dotted] (0,\y)--(14.5,\y);
\end{tikzpicture}

\newpage

%--- Problem 3: 25N.2.SL.TZ1.6 (P2 SL, calc) ---
\item Two machines are used in a factory to manufacture semiconductors.
      Machine A manufactures defective semiconductors $10\%$ of the time,
      while machine B manufactures defective semiconductors $5\%$ of the time.
      A randomly selected machine manufactures ten semiconductors.
      Both machines are equally likely to be selected.

      \begin{enumerate}[itemsep=0.1cm]
        \item Find the probability that exactly two semiconductors are defective. \hfill [4]
        \item Given that exactly two semiconductors are defective, find the
              probability that they were manufactured by machine A. \hfill [3]
      \end{enumerate}

\begin{tikzpicture}
  \draw (-0.4,0) rectangle (15,12);
  \foreach \y in {1,...,11} \draw[dotted] (0,\y)--(14.5,\y);
\end{tikzpicture}

\newpage

%--- Problem 4: 23M.2.SL.TZ1.5 (P2 SL, calc; uses z-score) ---
\item A company manufactures metal tubes for bicycle frames.
      The diameters of the tubes, $D$ mm, are normally distributed with
      mean $32$ and standard deviation $\sigma$. The interquartile range
      of the diameters is $0.28$.

      Find the value of $\sigma$. \hfill [5]

\begin{tikzpicture}
  \draw (-0.4,0) rectangle (15,12);
  \foreach \y in {1,...,11} \draw[dotted] (0,\y)--(14.5,\y);
\end{tikzpicture}

\newpage

%--- Problem 5: 25M.2.SL.TZ1.5 (P2 SL, calc; regression + normal+z) ---
\item In a study, measurements for arm span, $A$ cm, and foot length,
      $F$ cm, are taken from a large group of adults.

      For this group, the regression line of $F$ on $A$ is
      $F = 0.335A - 32.6$, and the regression line of $A$ on $F$ is
      $A = 2.89F + 99.3$. Each regression line passes through the mean point.

      \begin{enumerate}[itemsep=0.1cm]
        \item By using an appropriate regression line, find an estimate of the
              arm span for an adult with a foot length of $19.8$\,cm. \hfill [2]
        \item For this group of adults, find the mean arm span and the
              mean foot length. \hfill [3]
      \end{enumerate}

      The heights, $H$ cm, of adults in the group can be modelled by a normal
      distribution with mean $163$\,cm and standard deviation $\sigma$\,cm.
      It is found that $88\%$ of the group have a height between $153$\,cm
      and $173$\,cm.

      \begin{enumerate}[itemsep=0.1cm,resume]
        \setcounter{enumii}{2}
        \item Find the value of $\sigma$. \hfill [3]
      \end{enumerate}

\begin{tikzpicture}
  \draw (-0.4,0) rectangle (15,12);
  \foreach \y in {1,...,11} \draw[dotted] (0,\y)--(14.5,\y);
\end{tikzpicture}

\end{enumerate}
\end{document}
